\documentclass{article}

\usepackage[utf8]{inputenc}
\usepackage[spanish,es-noshorthands]{babel}
\usepackage{amssymb, amsmath, amsbsy}
\usepackage{tasks}
\usepackage{graphicx}

\fontfamily{cmss}\selectfont
\renewcommand{\familydefault}{\sfdefault}
\pagestyle{empty}

\begin{document}
% {author} -- Luis Carrillo Gutiérrez
\subsection*{Relaciones, dominio y rango}

\subsubsection*{Par ordenado}
\noindent Un {\tt par ordenado} de componentes ``\scalebox{1.2}{$\,a\,$}'' y ``\scalebox{1.2}{$\,b\,$}'' es un ente matemático denotado por $(a;\,b)$ donde ``$a$'' es la primera componente y ``$b$'' la segunda componente. Se define como : $(a;\,b) = \{\{a\};\,\{a;\,b\}\}$

\subsubsection*{Igualdad de pares ordenados}
\noindent Dos pares ordenados $(a;\,b)$ y $(c;\,d)$ son iguales sí y sólo si sus primeras componentes son iguales $a = c$, así como también sus segundas componentes $b = d$, en forma simultánea, es decir :
   
$$
(a;\,b) = (c;\,d) \iff a = c \wedge b = d
$$

\paragraph*{Ejemplo} -- Determinar el valor de \quad$x \wedge y$\quad y de tal manera que : \\ $(5x + 2y;\,-4) = (-1;\,2x - y)$

\begin{eqnarray*}
5x + 2y &=& -1 \\
2x - y &=& -4 \\
\text{Si}\quad y = 2x + 4 &\implies& 5x + 2 (2x + 4) = -1 \\
5x + 4x + 8 = -1 &\implies& x = -1\quad \wedge\quad y = 2
\end{eqnarray*}

\subsubsection*{Producto cartesiano}
\noindent Dados dos conjuntos $A$ y $B$ no vacíos se llama {\tt producto cartesiano} de $A$ y $B$ en ese orden al conjunto formado por todos los pares ordenados $(a;\,b)$ tal que $a \in A \wedge b \in B$ se denota por $A \times B$ esto es :
$$
A\times B = \{(a;\,b)/ a \in A \;\wedge\; b \in B\}
$$

\paragraph*{Ejemplo} -- Sean los conjuntos : \\ $A = \{1;\,2;\,3\}$\quad y \quad$B = \{2;\,4\}$, entonces \\ $A\times B = \{(1;\,2), (1;\,4), (2;\,2), (2;\,4), (3;\,2), (3;\,4)\}$

\paragraph*{Observación} -- Cuando los conjuntos $A$ y $B$ son finitos, y sean $p$ y $q$ la cantidad de elementos de cada conjunto respectivo, entonces el producto cartesiano $A\times B$ tiene $p\cdot q$ elementos, es decir :
$$ n(A\times B) = n(A) \cdot n(B) = p\cdot q$$
\paragraph*{Ejemplo} -- Sean $A = \{1;\,3;\,4\}$ y  $B = \{2;\,5;\,7;\,9\}$, entonces \\ $n(A\times B) = n(A) \cdot n(B) = 3 \cdot 4  = 12$\;pares ordenados.
\subsubsection*{Propiedades} % del producto cartesiano
\begin{itemize}
\item{Si $A \neq B$, entonces $A\times B \neq B\times A$}
\item{$A\times\varnothing = \varnothing\times A = \varnothing$}
\item{$A\times (B \times C) = (A \times B) \times C$}
\item{$A\times (B \cup C) = (A\times B) \cup (A\times C) $}
\item{$A\times (B \cap C) = (A\times B) \cap (A\times C)$}
\item{$A\times (B - C) = (A\times B) - (A\times C)$}
\item{Si $A \subset B$, entonces $(A\times C) \subset (B\times C)\;;\;\forall\; C$}
\end{itemize}
\end{document}